% !TEX root = ../main.tex
\section{Methods}
\label{sec:methods}

\subsection{Inventory binning}
\label{sec:methods_binning}

The legacy workflow operates on permanent sample plots drawn from the Qu\'ebec provincial inventory. For each species-group / cover-type combination we aggregated fixed-area plot tallies into 2~cm diameter classes spanning 10--60~cm DBH. The raw tallies $t_{ij}$ record the number of stems in plot $j$ falling in bin $i$. Because each plot covers $0.04$~ha, stems per hectare are obtained by multiplying by the constant expansion factor $s_{\text{plot}} = 1 / 0.04 = 25$. The empirically observed stand table for bin $i$ is therefore $\hat{y}_i = s_{\text{plot}} \, t_i$, where $t_i$ is the mean tally across plots.

Following the original notebook, we normalise binned observations so that the histogram integrates to one, i.e.\ $\sum_{i \in I} W \hat{y}_i^{(n)} = 1$ for bin width $W = 2$~cm. The vector $\bm{\hat{Y}} = \{\hat{y}_i^{(n)}\}$ represents a truncated empirical probability density function supported on $[a, b] = [10, 60]$~cm.

\subsection{Weighted non-linear least squares}
\label{sec:methods_nlls}

We fit parametric distributions by solving the weighted non-linear least squares problem
\begin{equation}
  Z = \min_{\bm{P}} \sum_{i \in I} e_i^2,
\end{equation}
with residuals
\begin{equation}
  e_i = w_i \left[f(x_i; \bm{P}) - \hat{y}_i^{(n)}\right],
\end{equation}
where $x_i$ denotes the class midpoint for bin $i$, $f$ is the probability density function evaluated at $x_i$, and $\bm{P}$ is the vector of free parameters. In the original legacy analysis, $w_i$ reflects the relative sampling error for bin $i$; because the public meta-plot data available to us no longer retain individual plot counts we set $w_i = 1$ for all bins.

Truncated datasets exhibit biased residual sums when complete-form densities are fit directly. The recommended approach rescales the target density to produce the truncated form
\begin{equation}
  f^{T}(x; \bm{P}) =
    \begin{cases}
      c\, f(x; \bm{P}) & a \leq x \leq b, \\
      0 & \text{otherwise},
    \end{cases}
\end{equation}
with scaling constant $c = [F(b; \bm{P}) - F(a; \bm{P})]^{-1}$ and $F$ the cumulative distribution function. When a closed form for $F$ is not available we evaluate the integral numerically.

\subsection{Two-stage fitting procedure}
\label{sec:methods_twostage}

We reproduce the two-stage complete-form workflow proposed in the legacy manuscript. For each candidate distribution we solve three variants:
\begin{enumerate}
  \item \textbf{1sc} -- a complete-form fit with the scaling fixed to $s = 1$;
  \item \textbf{1st} -- a truncated likelihood fit using $f^{T}(x; \bm{P})$ on $[a, b]$;
  \item \textbf{2sc} -- a two-stage complete-form fit that first estimates $s$ jointly with $\bm{P}$ and subsequently re-optimises with $s$ fixed at the stage-one estimate.
\end{enumerate}
We use the \texttt{lmfit} package to wrap each probability density function into a \texttt{Model} object, supply initial parameter guesses based on moments, and rely on the Levenberg--Marquardt optimiser to minimise $Z$.

\subsection{Candidate distributions and model selection}
\label{sec:methods_candidates}

The original analysis evaluated a broad family of generalised beta distributions. In this reimplementation we concentrate on the two distributions that dominated the published results: the two-parameter Weibull and gamma distributions (both specialisations of the generalised gamma). Each species-group / cover-type meta-plot is fit with both distributions under the 1sc, 1st, and 2sc regimes. We compute the small-sample Akaike information criterion,
\begin{equation}
  \text{AICc} = \text{AIC} + \frac{2K(K+1)}{n - K - 1},
\end{equation}
where $n$ is the number of non-zero bins and $K$ counts the number of free parameters plus the implicit error variance term. Consistent with the legacy manuscript we adjust $K$ to account for the additional scaling terms: $K_{1sc} = |\bm{P}| + 1$, $K_{1st} = |\bm{P}| + 3$ (two truncation bounds plus variance), and $K_{2sc} = |\bm{P}| + 2$. The best model for each meta-plot is selected by minimising AICc. We report parameter estimates, associated standard errors, and AICc values for both stage-one and stage-two fits to facilitate comparison with the truncated baseline.
